\section{Related Work}


\label{sec:related}
\noindent \textbf{Feed-Forward Novel-View Synthesis.} The shift from per-scene optimization to generalizable feed-forward networks has enabled 3DGS estimation without scene-specific training. Previous methods~\cite{charatan2024pixelsplat, chen2024mvsplat, xu2024depthsplat} explore pixel-aligned prediction by mapping 2D pixels to 3D Gaussians with transformers. While NoPoSplat~\cite{ye2024nopo}, Splatt3R~\cite{smart2024splatt3r}, and PF3plat~\cite{hong2024pf3plat} eliminated pose requirements, enabling zero-shot generalization from uncalibrated pairs. Recent models scale these ideas further, where FLARE~\cite{zhang2025flare} and AnySplat~\cite{jiang2025anysplat} unify sparse/dense uncalibrated collections, and VolSplat~\cite{wang2025volsplat} improved this with voxel-aligned prediction, directly regressing primitives from a 3D grid to avoid unreliable 2D matching. Despite these advances, all feed-forward methods fail in under-constrained regions where sparse observations yield geometric ambiguity.
\\
\\
% \paragraph{Sparse-View Novel-View Synthesis.} Sparse-view NVS requires hallucinating content from minimal observations. MVSplat~\cite{chen2024mvsplat} constructed plane-sweep cost volumes to establish cross-view correspondences, enabling reconstruction from two views. However, cost volumes struggle with extreme sparsity. MVSplat360~\cite{chen2024mvsplat360} addressed this by integrating feed-forward 3DGS with Stable Video Diffusion, rendering features into SVD's latent space for temporally consistent 360° synthesis from five views. PixelSplat~\cite{charatan2024pixelsplat} remains a calibrated baseline, while LEAP removes pose requirements via learned estimation. These methods reconstruct observed regions faithfully but cannot fill large missing areas while maintaining geometric consistency.
\noindent \textbf{Generation Control of Diffusion Model.} Controlling diffusion models beyond text prompts has become essential for spatial guidance. ControlNet~\cite{zhang2023ControlNet} explores trainable zero-convolutions for injecting structural conditions, while T2I-Adapter~\cite{mou2024t2i} provides a lighter-weight composable alternative. Dhariwal \etal~\cite{dhariwal2021diffusion} provides an additional noise-conditioned external classifier guidance for the denosing sampler. Training-free approaches like Universal Guidance~\cite{bansal2023universal} and FREEDOM~\cite{yu2023freedom} avoid per-condition training by steering generation via gradients from pre-trained models. However, all these methods operate purely in the 2D latent space and lack explicit geometric reasoning, ill-suited for maintaining multi-view consistency in 3D reconstruction. Our geometry-aware leveling adapter bridges this gap by directly conditioning on multi-view geometry prior, aligning diffusion priors with feed-forward 3D geometry to enable simultaneous reconstruction and generation.
\\
\\
\noindent \textbf{3D Reconstruction Enhancement with Generation.} Recent work has used diffusion models as post-processors to repair reconstruction artifacts. Difix3D+~\cite{wu2025difix3d+}  trained a single-step diffusion model on artifact-clean pairs and distilled priors into 3DGS optimization. GSFixer~\cite{yin2025gsfixer} advanced this with a DiT-based architecture that fuses 2D semantic and 3D geometric features from reference views to ensure multi-view consistency. GSFix3D~\cite{wei2025gsfix3d} similarly distilled diffusion priors using random mask augmentation. Moreover, Zhong \etal~\cite{zhong2025taming} used video diffusion to iteratively enhance views with the training-free guidance approach. However, these methods solely rely on image-based reference and operate as post-processing tools, failing at improving the underlying geometry representation. Leveling3D departs by tightly coupling reconstruction with geometry-aware diffusion in a unified framework.